% !PRINTED 702
\weberpage{702}
%<<CONT
associated integers divisible by $\varphi_1$ whose norms are not greater than $t_1$, and by the theorem \S.\,190,\,1. we have
\weq{2}{\mathop{\mathrm{Lim}}\limits_{t_1=\infty}\frac{T_1}{t_1} = \frac{g}{N_a(\varphi_1)},}
where $g$ has the meaning given at the place referred to, and is accordingly a number \gesperrt{different from zero}, determined by the nature of the field $\Omega$. If now $T_2, t_2, \ldots, T_h, t_h$ have the corresponding meaning for the classes $A_2, \ldots, A_h$ that $T_1, t_1$ have for $A_1$, then
\[
T = T_1 + T_2 + \cdots + T_h,
\]
\[
\frac{T}{t} = \frac{T_1}{t_1}\, N_a(\varphi_1) + \frac{T_2}{t_2}\, N_a(\varphi_2) + \cdots + \frac{T_h}{t_h}\, N_a(\varphi_h),
\]
and from (1) and (2) the fundamental formula results:
\weq{3}{\mathop{\mathrm{Lim}}\limits_{s=1}\sum \frac{s-1}{N(\mathfrak{a})^s} = gh.}
% [sic] The following source citation is section 185,(11); the formula defining g is section 190,(11) on printed p. 693.
The number $g$ we may regard as known by \S.\,185,\,(11), even though its actual calculation still rests on the assumption that a fundamental system of units is known, and even though the ascertaining of such a system in higher fields has always been regarded as one of the greatest difficulties. The calculation of the class number then still depends on the determination of the limiting value on the left side of (3), which of course likewise succeeds only in special cases, but which often permits important conclusions about the nature of the class number. We wish to develop one further transformation, preparatory to the calculation, of the sum
\weq{4}{\sum \frac{1}{N(\mathfrak{a})^s} = \Phi(s)}
into an infinite product.

Let $\mathfrak{p}$ be any prime ideal of degree $f$ in the field $\Omega$, so that
\[
N(\mathfrak{p}) = p^f.
\]
Then the infinite geometric series is
\[
1 + \frac{1}{N(\mathfrak{p})^s} + \frac{1}{N(\mathfrak{p})^{2s}} + \cdots = \frac{1}{1-p^{-sf}};
\]
if one multiplies these series for all the different prime ideals $\mathfrak{p}$ with one another, then by known theorems one obtains from
%CONT>>
